\section{Introduction}
\label{sec:intro}

The carbon intensity of electricity differs widely between grids and
over the day. In January 2024 the Norwegian grid averaged about
28\,gCO$_2$eq/kWh while Poland's averaged 586, and Germany's ranged from
130 to 650 within days as wind and solar output changed
(\S\ref{sec:method}). A cloud job emits very different amounts of
carbon depending on \emph{where} and \emph{when} it runs. Carbon-aware
schedulers exploit this through \emph{spatial shifting} (moving work to
cleaner regions)~\cite{zhou2013carbon,gsteiger2024caribou} and
\emph{temporal shifting} (delaying work to cleaner
hours)~\cite{wiesner2021wait,radovanovic2023carbon}.

Acun et al.~\cite{acun2024limitations} showed that the realistic gains
of both forms of shifting are smaller than idealised analyses suggest.
Much of the benefit is captured by simple policies, and capacity,
latency and data movement limit the rest. Three practical questions
follow:
\begin{itemize}
\item What does shifting cost in \emph{service quality}, that is, how
  many deadlines does it break on a real, bursty workload?
\item How much does the cost of moving data shrink spatial shifting?
\item How robust are forecast-driven policies when forecasts are wrong?
\end{itemize}
Published policies are rarely compared on these questions under the
same constraints.

This paper studies these questions with \emph{GreenOps}, a
multi-objective scheduler that decides, for every job, a region and a
start time. It weighs forecast carbon (including data-transfer
emissions), cost and a convex measure of deadline risk. Deadlines,
budgets and data residency are hard constraints. We evaluate GreenOps
against six baselines in a discrete-event simulator. The baselines
include the temporal-shifting strategy of Wiesner et
al.~\cite{wiesner2021wait} and a forecast-free threshold rule underlying
the competitive algorithms of Lechowicz et al.~\cite{lechowicz2024soad}.
The simulator is driven by real 2024 grid data and $\approx$45\,k jobs
per run from the Alibaba 2018 cluster trace~\cite{guo2019alibaba}.

Our findings are:
\begin{itemize}
\item \textbf{Deferral protects deadlines.} On the real trace, policies
  that never delay work miss the deadlines of about 5\,\% of urgent jobs
  at demand peaks. Delaying flexible work frees capacity, and GreenOps
  has the fewest violations of all policies.
\item \textbf{There is no free lunch.} A forecast-free threshold policy
  reaches the lowest carbon, but at several times GreenOps' violation
  rate. No GreenOps weighting matches its carbon. The policies trace a
  genuine carbon--SLA frontier rather than a single winner.
\item \textbf{Data gravity acts mainly through cost.} At realistic data
  volumes, network energy is a small share of emissions. Egress fees
  and transfer time, not transfer carbon, are what pull jobs back to
  their home region.
\item \textbf{Re-planning beats trusting the forecast.} Policies that
  commit to a forecast once degrade fastest as forecast error grows.
\end{itemize}
The simulator, traces and experiment scripts are open
source.\footnote{\url{https://github.com/SANJEEVKUMARKS-24/greenops-cloud-orchestration}}
